% GENERATED FILE. Edit canonical lessons, then run npm run generate:books.
% canonical-source-hash: fnv1a64:69e3cda081365e01
% canonical-lessons: UR-C144-hal-hi-mein, UR-C144-pehle-se, UR-C144-mazi, UR-C144-guzashta, UR-C144-pesh-ana

\chapter{Recently, Already, the Past, Last, To Happen}
\label{ch:ur-recently-already-the-past-last-to-happen}
\hlchaptermodality{144}

\begin{chapteropening}
\textbf{By the end of this chapter:} \emph{I can say what happened, and when: recently, already, last year.}

Complete the last lesson of chapter 144: I can say what happened, and when: recently, already, last year.
\end{chapteropening}

\section[\texorpdfstring{\ur{حال} \ur{ہی} \ur{میں} (hāl hī meṅ)}{حال ہی میں (hāl hī meṅ)}]{\ur{حال} \ur{ہی} \ur{میں} (hāl hī meṅ) — recently}

\label{lesson:UR-C144-hal-hi-mein}



\begin{quote}
Before the new one: say the Urdu for to be saved, to avoid, then the Urdu for to grow, to increase.
\end{quote}

\subsection*{You'll want to know: \ur{حال} \ur{ہی} \ur{میں}}
\textbf{\ur{حال} \ur{ہی} \ur{میں}} — \emph{hāl hī meṅ} — \textquotedblleft{}recently\textquotedblright{}. With a verb like \textbf{\ur{پڑھنا}}, Urdu's past puts \textbf{\ur{نے}} after the doer: \textbf{\ur{میں} \ur{نے} \ur{حال} \ur{ہی} \ur{میں} \ur{ایک} \ur{کتاب} \ur{پڑھی۔}} — \emph{maiṅ ne hāl hī meṅ ek kitāb paṛhī} — \textquotedblleft{}I read a book recently.\textquotedblright{} \textbf{\ur{پڑھی}} agrees with \textbf{\ur{کتاب}}.

\begin{cousinweb}
\textbf{\ur{حال} \ur{ہی} \ur{میں}} (\emph{hāl hī meṅ}) is \textquotedblleft{}in the very present state\textquotedblright{}. \emph{hāl} is a loan from Arabic \emph{ḥāl}, \textquotedblleft{}state, condition\textquotedblright{}, through Persian (Lane, root \emph{ḥ-w-l}); \emph{hī} is the emphatic \textquotedblleft{}very\textquotedblright{}; and \emph{meṅ}, \textquotedblleft{}in\textquotedblright{}, comes from Sanskrit \emph{madhye}, \textquotedblleft{}in the middle\textquotedblright{} (CDIAL 9804).
\end{cousinweb}

\begin{grammarlens}[title={What you've built}]
One more word for this chapter.
\end{grammarlens}

\subsection*{Your turn}
\begin{itemize}
\raggedright
  \item \emph{Say it:} \emph{hāl hī meṅ}
  \item \emph{Say it:} \emph{hāl hī meṅ}, once more
  \item \emph{Say it:} say \emph{hāl hī meṅ}
  \item \emph{Recall:} say the Urdu for to be saved, to avoid, then the Urdu for to grow, to increase, then say \emph{hāl hī meṅ} again
\end{itemize}

\begin{tcolorbox}[breakable,colback=teal!4,colframe=teal!35!black,title={Before you move on}]
What does \ur{حال} \ur{ہی} \ur{میں} mean? (\textquotedblleft{}Recently\textquotedblright{}.) Say it once more.
\end{tcolorbox}
\section[\texorpdfstring{\ur{پہلے} \ur{سے} (pehle se)}{پہلے سے (pehle se)}]{\ur{پہلے} \ur{سے} (pehle se) — already, beforehand}

\label{lesson:UR-C144-pehle-se}



\begin{quote}
Before the new one: say the Urdu for to grow, to increase, then the Urdu for recently.
\end{quote}

\subsection*{You'll want to know: \ur{پہلے} \ur{سے}}
\textbf{\ur{پہلے} \ur{سے}} — \emph{pehle se} — \textquotedblleft{}already, beforehand\textquotedblright{}. It is \textbf{\ur{پہلے}}, \textquotedblleft{}before\textquotedblright{}, with \textbf{\ur{سے}}: \textbf{\ur{میں} \ur{پہلے} \ur{سے} \ur{جانتا} \ur{ہوں۔}} — \emph{maiṅ pehle se jāntā hūṅ} — \textquotedblleft{}I already know.\textquotedblright{} A woman says \textbf{\ur{جانتی}}.

\begin{cousinweb}
\textbf{\ur{پہلے} \ur{سے}} (\emph{pehle se}) joins \emph{pehle}, \textquotedblleft{}first, before\textquotedblright{}, and \emph{se}, \textquotedblleft{}from\textquotedblright{}. \emph{pehle} is the adverb of \emph{pahlā}, which CDIAL lists under \emph{*prathilla}, a variant of a reconstructed \emph{*prathila}, \textquotedblleft{}first\textquotedblright{} (CDIAL 8652). CDIAL lists \emph{se} under Sanskrit \emph{sahita}, \textquotedblleft{}together with\textquotedblright{} (CDIAL 13310), and \emph{sama}, \textquotedblleft{}equal\textquotedblright{}, has also been suggested (CDIAL 13173).
\end{cousinweb}

\begin{grammarlens}[title={What you've built}]
One more word for this chapter.
\end{grammarlens}

\subsection*{Your turn}
\begin{itemize}
\raggedright
  \item \emph{Say it:} \emph{pehle se}
  \item \emph{Say it:} \emph{pehle se}, once more
  \item \emph{Say it:} say \emph{pehle se}
  \item \emph{Recall:} say the Urdu for to grow, to increase, then the Urdu for recently, then say \emph{pehle se} again
\end{itemize}

\begin{tcolorbox}[breakable,colback=teal!4,colframe=teal!35!black,title={Before you move on}]
What does \ur{پہلے} \ur{سے} mean? (\textquotedblleft{}Already, beforehand\textquotedblright{}.) Say it once more.
\end{tcolorbox}
\section[\texorpdfstring{\ur{ماضی} (māzī)}{ماضی (māzī)}]{\ur{ماضی} (māzī) — the past}

\label{lesson:UR-C144-mazi}



\begin{quote}
Before the new one: say the Urdu for recently, then the Urdu for already.
\end{quote}

\subsection*{You'll want to know: \ur{ماضی}}
\textbf{\ur{ماضی}} — \emph{māzī} — \textquotedblleft{}the past\textquotedblright{}. Grammar books use it for the past tense too. It came into Urdu from Arabic.

\begin{cousinweb}
\textbf{\ur{ماضی}} (\emph{māzī}) is a loan from Arabic \emph{māḍī}, \textquotedblleft{}past, gone by\textquotedblright{}, which reached Urdu through Persian. Its Arabic root is \emph{m-ḍ-y} (Lane). Grammar books use it for the past tense.
\end{cousinweb}

\begin{grammarlens}[title={What you've built}]
One more word for this chapter.
\end{grammarlens}

\subsection*{Your turn}
\begin{itemize}
\raggedright
  \item \emph{Say it:} \emph{māzī}
  \item \emph{Say it:} \emph{māzī}, once more
  \item \emph{Say it:} say \emph{māzī}
  \item \emph{Recall:} say the Urdu for recently, then the Urdu for already, then say \emph{māzī} again
\end{itemize}

\begin{tcolorbox}[breakable,colback=teal!4,colframe=teal!35!black,title={Before you move on}]
What does \ur{ماضی} mean? (\textquotedblleft{}The past\textquotedblright{}.) Say it once more.
\end{tcolorbox}
\section[\texorpdfstring{\ur{گزشتہ} (guzashtā)}{گزشتہ (guzashtā)}]{\ur{گزشتہ} (guzashtā) — last, previous}

\label{lesson:UR-C144-guzashta}



\begin{quote}
Before the new one: say the Urdu for already, then the Urdu for the past.
\end{quote}

\subsection*{You'll want to know: \ur{گزشتہ}}
\textbf{\ur{گزشتہ}} — \emph{guzashtā} — \textquotedblleft{}last, previous\textquotedblright{}. \textbf{\ur{گزشتہ} \ur{سال}} — \emph{guzashtā sāl} — \textquotedblleft{}last year\textquotedblright{}. It is the formal word; everyday speech also says \textbf{\ur{پچھلا}}.

\begin{cousinweb}
\textbf{\ur{گزشتہ}} (\emph{guzashtā}) is a loan, usually traced to Persian \emph{guzashta}, \textquotedblleft{}passed, gone by\textquotedblright{}. It is the past participle of \emph{guzashtan}, \textquotedblleft{}to pass\textquotedblright{}.
\end{cousinweb}

\begin{grammarlens}[title={What you've built}]
One more word for this chapter.
\end{grammarlens}

\subsection*{Your turn}
\begin{itemize}
\raggedright
  \item \emph{Say it:} \emph{guzashtā}
  \item \emph{Say it:} \emph{guzashtā}, once more
  \item \emph{Say it:} say \emph{guzashtā}
  \item \emph{Recall:} say the Urdu for already, then the Urdu for the past, then say \emph{guzashtā} again
\end{itemize}

\begin{tcolorbox}[breakable,colback=teal!4,colframe=teal!35!black,title={Before you move on}]
What does \ur{گزشتہ} mean? (\textquotedblleft{}Last, previous\textquotedblright{}.) Say it once more.
\end{tcolorbox}
\section[\texorpdfstring{\ur{پیش} \ur{آنا} (pesh ānā)}{پیش آنا (pesh ānā)}]{\ur{پیش} \ur{آنا} (pesh ānā) — to happen, to occur}

\label{lesson:UR-C144-pesh-ana}



\begin{quote}
Before the new one: say the Urdu for the past, then the Urdu for last.
\end{quote}

\subsection*{You'll want to know: \ur{پیش} \ur{آنا}}
\textbf{\ur{پیش} \ur{آنا}} — \emph{pesh ānā} — \textquotedblleft{}to happen, to occur\textquotedblright{}. Said of events: \textbf{\ur{کل} \ur{ایک} \ur{مسئلہ} \ur{پیش} \ur{آیا۔}} — \emph{kal ek mas'alā pesh āyā} — \textquotedblleft{}A problem came up yesterday.\textquotedblright{} Its past is \textbf{\ur{پیش} \ur{آیا}}, or \textbf{\ur{پیش} \ur{آئی}} with a feminine subject.

\begin{cousinweb}
\textbf{\ur{پیش} \ur{آنا}} (\emph{pesh ānā}) is \textquotedblleft{}to come to the front\textquotedblright{}. \emph{pesh}, \textquotedblleft{}front, before\textquotedblright{}, is usually traced to Persian \emph{pesh}; \emph{ānā}, \textquotedblleft{}to come\textquotedblright{}, comes from Sanskrit \emph{āpayati}, \textquotedblleft{}causes to reach\textquotedblright{} (CDIAL 1200).
\end{cousinweb}

\begin{grammarlens}[title={What you've built}]
One more word for this chapter.
\end{grammarlens}

\subsection*{Your turn}
\begin{itemize}
\raggedright
  \item \emph{Say it:} \emph{pesh ānā}
  \item \emph{Say it:} \emph{pesh ānā}, once more
  \item \emph{Say it:} say \emph{pesh ānā}
  \item \emph{Recall:} say the Urdu for the past, then the Urdu for last, then say \emph{pesh ānā} again
\end{itemize}

\begin{tcolorbox}[breakable,colback=teal!4,colframe=teal!35!black,title={Before you move on}]
What does \ur{پیش} \ur{آنا} mean? (\textquotedblleft{}To happen, to occur\textquotedblright{}.) Say it once more.
\end{tcolorbox}
